% from ega0-19-fr.tex
\oldpage[0\textsubscript{IV-1}]{83}

\begin{proposition}[19.3.6]
\label{0.19.3.6-fr}
Soient $A$ un anneau topologique, $B$ une $A$-algèbre topologique, $\hat A$
et $\hat B$ les séparés complétés respectifs de $A$ et $B$, de sorte que
$\hat B$ est une $\hat A$-algèbre topologique. Les conditions suivantes
sont équivalentes~:
\begin{enumerate}
  \item[a)] $B$ est une $A$-algèbre formellement lisse.
  \item[b)] $\hat B$ est une $A$-algèbre formellement lisse.
  \item[c)] $\hat B$ est une $\hat A$-algèbre formellement lisse.
\end{enumerate}
\end{proposition}

\begin{proof}
Bien entendu, la structure de $\hat A$-algèbre sur $\hat B$ est définie
par l'homomorphisme $\hat\varphi$, si $\varphi:A\to B$ est
l'homomorphisme définissant la structure de $A$-algèbre sur $B$. Comme toute
$A$-algèbre discrète $C$ est séparée et complète, elle est aussi une
$\hat A$-algèbre (par prolongement canonique de l'homomorphisme de $A$ dans
$C$), et un $A$-homomorphisme continu de $B$ dans $C$ donne par prolongement
canonique un $\hat A$-homomorphisme (et \emph{a fortiori} un
$A$-homomorphisme) de $\hat B$ dans $C$ (autrement dit, tout
$A$-homomorphisme $B\to C$ se factorise en $B\to\hat B\to C$ de façon
unique). Ces remarques et la définition (19.3.1) entraînent aussitôt la
proposition.
\end{proof}

\begin{corollary}[19.3.7]
\label{0.19.3.7-fr}
Sous les hypothèses de (19.3.5), (iv), la $A\{S^{-1}\}$-algèbre topologique
$B\{T^{-1}\}$ est formellement lisse.
\end{corollary}

Cela résulte des définitions ($\mathbf 0_{\rm I}$, 7.6.1 et 7.6.7) et de
(19.3.5), (iv) et (19.3.6).

\begin{proposition}[19.3.8]
\label{0.19.3.8-fr}
Soient $A$ un anneau topologique, $B$ une $A$-algèbre topologique, et supposons
que pour tout idéal ouvert $\mathfrak R$ de $B$, $\mathfrak R^2$ soit aussi
ouvert. Soient $A'$, $B'$ des anneaux topologiques obtenus en munissant $A$ et
$B$ de topologies plus fines que les topologies initiales, et supposons que
l'homomorphisme canonique $\varphi:A\to B$ soit encore un homomorphisme
continu de $A'$ dans $B'$. Alors, si $B'$ est une $A'$-algèbre formellement
lisse, $B$ est une $A$-algèbre formellement lisse.
\end{proposition}

Il suffit d'appliquer le lemme suivant~:

\begin{lemma}[19.3.8.1]
\label{0.19.3.8.1-fr}
Soient $C$ une $A$-algèbre discrète, $\mathfrak J$ un idéal nilpotent de $C$.
Supposons que le carré de tout idéal ouvert de $B$ soit ouvert. Alors, si un
homomorphisme composé
$v:B\xrightarrow{u}C\to C/\mathfrak J$ est continu, l'homomorphisme $u$
est continu.
\end{lemma}

\begin{proof}
En effet, $\mathfrak R=u^{-1}(\mathfrak J)$ est le noyau de $v$ et il est
ouvert par hypothèse~; comme il existe un entier $n$ tel que
$\mathfrak J^n=0$, $\mathfrak R^n$ est contenu dans
$\operatorname{Ker}(u)$~; mais par récurrence sur $h$, il résulte de
l'hypothèse que $\mathfrak R^{2^h}$ est ouvert, donc aussi
$\mathfrak R^n$ pour tout $n$, et par suite $\operatorname{Ker}(u)$ est lui
aussi ouvert, ce qui prouve notre assertion.
\end{proof}

On notera que l'hypothèse sur $B$ signifie que la topologie de $B$ est borne
supérieure des topologies $\mathfrak R$-préadiques, où $\mathfrak R$ parcourt
l'ensemble des idéaux ouverts de $B$. Si $B$ est un anneau préadique
($\mathbf 0_{\rm I}$, 7.1.9), cette condition est toujours vérifiée.

\begin{env}[19.3.9]
\label{0.19.3.9-fr}
Nous allons maintenant voir que la propriété d'être formellement lisse
implique des propriétés de «~relèvement~» d'homomorphismes sous des conditions
plus générales que celles de la définition (19.3.1).
\end{env}

